\chapter*{Introduction}\label{ch:introduction}
In recent years, companies' rapid digital transformation has influenced the cybersecurity landscape, mainly because they have decided to move their assets to remote servers, typically in the cloud, by exploiting technologies like cloud computing and the Internet of Things (IoT). Remote work positions have expanded the attack surface as employees access corporate networks from different locations and devices.

In this scenario, companies must update their cybersecurity strategies to limit the risks of data breaches and other security attacks. According to the study of Dave et al.~\cite{dave2023}, cybersecurity is a growing concern due to the rapid advancement of technology and the increasing interconnectedness of systems. As main threats, there are malware attacks (e.g., viruses, worms, Trojans, ransomware, spyware, adware), social engineering attacks (e.g., phishing, pretexting, tailgating), network vulnerabilities, and data breaches.

One of the trending cyberattacks is ransomware, where malicious software encrypts the victim's data and makes it inaccessible until the victims pay a ransom. During the COVID-19 pandemic months, organizations started to shift toward remote work positions because of the law restrictions. Still, it highlighted new vulnerabilities in the company systems since the employees had to access the work programs from different points in the network, often with their devices.

Regardless of the threat type, there are increasing attacks on IoT devices like smart appliances, wearable technologies, home automation systems, and cloud-based services, such as compromising cloud storage accounts or stealing sensitive data from cloud-based applications. In light of these developments, awareness and vigilance have become a high priority: ``By remaining informed and vigilant, we can collectively work toward forging a safer and more secure digital future.''~\cite{dave2023}

In response to this expanding landscape, the European Union funded the CS-Aware-Next project as part of the Horizon Europe program. This project builds on the results obtained during the previous initiative, CS-Aware, which concluded in August 2020 and provided an advanced platform for the awareness and management of cybersecurity in real time for local government and Network and Information Security Directive (NIS) sector organizations. CS-Aware-Next offers cybersecurity management policies for organizations to face the dynamic, ever-evolving cybersecurity landscape. Moreover, the project improves threat intelligence integration in cybersecurity management by exploiting innovative artificial intelligence approaches and techniques.

Building on these developments, the project helps organizations conform to European cybersecurity laws, such as the NIS, which sets the requirements organizations must meet to guarantee a high level of cybersecurity. For this reason, CS-Aware-Next also integrates classic risk and incident management actions to offer dynamic emergency recovery.

The European project recognizes the crucial role of social media in managing cybersecurity. On the one hand, social platforms are essential tools for communication and the spread of information, allowing organizations to share updates on urgent threats and vulnerabilities rapidly. On the other hand, they represent a significant challenge because they are repositories of vast amounts of content, much of which may be irrelevant, misleading, or even malicious. This creates a pressing need to filter and identify relevant cybersecurity information.

Amongst the objectives of CS-Aware-Next that align with the goals of this thesis are dynamic cybersecurity management, regional collaboration and cooperation, and threat intelligence integration. Regarding dynamic cybersecurity management, the project aims to improve cybersecurity by increasing small companies' awareness and encouraging collaboration. At the same time, the thesis contributes by developing a model to filter and categorize cybersecurity-related social media posts, which integrates threat intelligence from various sources, including social media, into organizational practices. This integration facilitates real-time risk and incident management~\cite{dionisio2019} and supports the development of a cybersecurity management policy framework for the evolving landscape. Furthermore, the model advances regional collaboration and cooperation by enabling real-time threat detection and information sharing—capabilities that enhance situational awareness and support cooperation between stakeholders at local, regional, and national levels. Finally, threat intelligence integration is furthered as the thesis develops methodologies and tools to automate and improve intelligence sharing, leveraging social media's key role in rapidly disseminating news, updates, and security notes.

This thesis looks at several important areas in cybersecurity. It starts with Cyber Threat Intelligence, which is about gathering and studying data from social media posts, blogs, and forums on the internet. The purpose is to find new cyber threats and trends and keep track of them. This information is used to create innovative ways to sort and filter text to pick out cybersecurity-related content from large amounts of information online, making it easier to organize. The study also deals with multi-label classification, where one piece of information can simultaneously fit into more than one category. This is challenging because the categories are linked, the data can be uneven, and careful handling of cybersecurity tasks is crucial. The thesis uses advanced word tools like GloVe and BERT to tackle these language issues. These tools are excellent for understanding the meaning and context of text, which is essential for sorting tasks. Finally, the research includes different algorithm methods. It uses basic machine learning models like logistic regression and random forests. It also explores more advanced techniques that combine several models to boost performance and accuracy in identifying and sorting cybersecurity threats.

Two system architectures were developed and compared to achieve these objectives. The Two-Step System separates the process into a filtering stage that discards non-relevant content and a categorization stage that assigns cybersecurity labels to the relevant social media posts. The One-Step System integrates these tasks into a single categorization model, in which the relevance of a post is implied by the prediction of at least one category. Both systems were built on a single grouped train-test split, with selection and tuning confined to the training partition, and were evaluated once on the test partition with group-bootstrap confidence intervals. On the primary endpoint, the micro F1 of the public labels, the two architectures are not distinguishable, while they differ in the balance between precision and recall. The embedder of each component was chosen by cross-validation, and the best one varies across components.

\section*{Roadmap}\label{sec:roadmap}
Here is an outline of the chapters presented in this work, each with a brief.
\begin{itemize}
    \item Chapter~\ref{ch:state-of-the-art} reviews the state of the art in cybersecurity analysis and text classification methodologies.
    \item Chapter~\ref{ch:methodology} details the methodology, including the datasets, preprocessing, taxonomy, model architectures, and evaluation protocol.
    \item Chapter~\ref{ch:experimental-results} presents the experimental results: the filter, the categorization levels, and the end-to-end comparison of the two systems.
    \item Chapter~\ref{ch:conclusions} finally concludes with a discussion of the findings, limitations, and directions for future research.
\end{itemize}